% Part: lambda-calculus
% Chapter: church-rosser
% Section: beta-eta-reduction

\documentclass[../../../include/open-logic-section]{subfiles}

\begin{document}

\olfileid{lam}{cr}{be}

\olsection{$\beta\eta$-न्यूनीकरण}

$\beta\eta$-न्यूनीकरणाला ($\bered$) चर्च–रॉसर
गुणधर्म आहे.

\begin{lem}\ollabel{lem:one-par}
  $M \beredone M'$ असेल, तर $M \beredpar M'$.
\end{lem}

\begin{proof}
  $M \beredone M'$ च्या !!{derivation} वर विगमन करू.
  $M \eredone M'$ हे $\eta$-परिवर्तनाचे प्रकरण
  असेल (म्हणजे \olref[syn][eta]{defn:beredone}),
  तर समांतर इटा-नियमासाठी \olref[pbe]{thm:refl}
  वापरू. इतर प्रकरणे \olref[b]{lem:one-par}
  प्रमाणेच आहेत.
\end{proof}

\begin{lem}\ollabel{lem:par-red}
  $M \beredpar M'$ असेल, तर $M \bered M'$.
\end{lem}

\begin{proof} $M \beredpar M'$ च्या !!{derivation} वर विगमन करू.

  पहिली चार प्रकरणे समांतर बीटा-न्यूनीकरणाच्या
  संबंधित प्रकरणांप्रमाणे हाताळता येतात. अखेरचा नियम
  \olref[pbe]{defn:beredpar5} असेल, तर काही $x$, $N$,
  $N'$ साठी $M$ हे $\lambd[x][Nx]$ आणि $M'$ हे $N'$
  आहे, जेथे $x \notin FV(N)$ आणि $N \beredpar N'$.
  म्हणून आधी $\lambd[x][Nx]$ चे $N$ मध्ये
  $\eta$-परिवर्तनाने न्यूनीकरण करू आणि मग विगमन
  गृहीतकानुसार $N \bered N'$ दाखवणाऱ्या $\beredone$
  पायऱ्या घेऊ.
\end{proof}

\begin{lem}\ollabel{lem:str}
  $\beredpar$ समाविष्ट करणारा सर्वांत लहान संक्रमक
  संबंध $\bered$ आहे.
\end{lem}

\begin{proof}
  \olref[b]{lem:str} प्रमाणे.
\end{proof}

\begin{thm}\ollabel{thm:cr}
  $\bered$ ला चर्च–रॉसर गुणधर्म आहे.
\end{thm}

\begin{proof}
  \olref[dap]{thm:str}, \olref[pbe]{thm:cr} आणि
  \olref{lem:str} वरून मिळते.
\end{proof}
\end{document}
